\subsection{从江都到长安：同一条水路不能保住同一个朝廷}
618年三月，骁果军在江都发动兵变，\eventref{ST-618-JIANGDU-COUP}{炀帝被杀}。这件事的顺序可靠，遗诏、责任和宫廷细节却被唐初的亡国叙事反复加工。与其追逐那些难以核实的细节，不如停在江都的船、仓和军营之间：一支长期驻在南方的军队为何要北返，沿河的粮食、船只和官吏又将听谁的命令。运河原本把江南、洛阳和关中接进同一调度网；政变发生后，水路仍在，能够调度它的中心却裂开了。

对沿线州县而言，军队北返不是抽象的改朝。船户会被征用，仓吏要决定粮食交给谁，军属和商人则要猜测下一处城门是否还会承认原来的文书。我们没有连续的地方账簿，不能替每个家庭写出同一种恐惧或选择；但正因为材料只在军队移动和城池易手时偶尔照见他们，才更应把“江都兵变”放进交通与供给的时间里。隋末的断裂不是运河失去作用，而是它继续把兵、人和消息更快送向互相竞争的中心。

五月，\eventref{ST-618-TANG-FOUNDING}{唐建国}的受禅仪式给关中一个新名义。仪式不能把江都、洛阳、河北和江南已发生的破裂收回。李渊首先接手的是不完整的仓簿、会变动的军队和需要逐地恢复的道路。唐能继承隋的行政与运输遗产，正是因为仍有人愿意守仓、驾船、传递公文；唐也必须用较低的立即索取来换取这种合作。隋唐之间的过渡因此不是一扇宫门关闭、另一扇开启，而是许多地点在不同月份重新决定谁的凭证可以换到粮食和保护。
